%% ============================================================================
\section{Discussion}
\label{sec:discussion}
%% ============================================================================

\subsection{AE vs. US}
\label{subsec:disc_ae}

The consistently superior performance of AE over the US channel across all model
families reflects a fundamental hardware difference. The AE sensor captures wideband, rotation-generated content extending into the MHz range (analysed here up to \SI{1}{MHz}), dominated by verified stress-wave emission up to several hundred kHz. The US channel down-mixes its $\approx$38--40~kHz sensing band to an audio-frequency baseband (Section~\ref{subsec:origins}). Although the probe acquires content across its 0--100~kHz band (Section~\ref{subsec:filter_bands}), the lubrication-relevant information is concentrated in the nominal 0--20~kHz baseband---a mixture of down-mixed contact activity, low-frequency machinery noise, VFD harmonics, and structural vibration, within which the 10--20~kHz half carries the strongest $\kappa$ correlations (Section~\ref{subsec:raw_spearman})---while the 20--100~kHz above-baseband content contributes comparatively little. The high-frequency transient signature of individual microscale contact events, which the AE channel captures directly, is largely absent from this down-mixed envelope. Replacing the heterodyned probe with a wideband passive acoustic or ultrasound sensor would be expected to substantially improve the US-channel models.

%% Former \S "Origin of the high-frequency AE content" (subsec:disc_hf) removed
%% 2026-06-15: obsolete after exclusion of the >1 MHz (1--2 MHz) band, on which the
%% whole subsection (mechanical vs. bearing-current/EDM hypotheses) was predicated.
%% NOTE: the still-valid deployment point that the 20--500 kHz band rises
%% monotonically with speed and should transfer across drive configurations was
%% lost with it -- relocate to a retained-band discussion if wanted.

\subsection{Fusion gain}
\label{subsec:disc_fusion}

The fusion gain ($\Delta R^2 = \resDrsqLgbCombAe$, $\Delta\text{RMSE} = \resDrmseAeComb$ on the hold-out) is small but statistically robust, detected under repeated nested cross-validation and confirmed by all three hold-out tests. Its limited magnitude can be understood from two complementary angles.
First, the US channel retains only \resNretUs{} features after selection (vs.\ \resNretAe{} for AE), and
the top US correlations ($|\rho| \leq \resTopUsRho$) are weaker than the top AE
correlations ($|\rho| = \resTopAeRho$). The US channel therefore contributes a genuine but limited amount of non-redundant information to the combined model. Second, the US channel's down-mixed audio-frequency baseband means its information content partially overlaps with the low-frequency tail of the AE signal, limiting the independent contribution it can make.

This result is an informative contribution to the LCM literature, where AE--ultrasound complementarity is widely \emph{argued} but rarely \emph{tested}.
Whether a wideband passive ultrasound sensor (rather than a heterodyned probe) would
provide more independent information -- and thus a more substantial fusion gain --
remains an open question for future work.

\subsection{Interpreting the amplitude-confounding result}
\label{subsec:disc_confound}

The strong raw Spearman correlations of amplitude features (RMS and related
statistics) and frequency-tracking features (Hjorth mobility, centre frequency)
with $\kappa$ should be interpreted with caution. As discussed in Section~\ref{subsec:confound}, rotational speed is a driver of both $\kappa$ (through $\nureq$ in Eq.~\eqref{eq:nu1}) and broadband signal energy (through the rate of and energy in contact events per revolution as described in~\cite{zhang_coupled_2026}). A model trained on amplitude features therefore exploits speed as an indirect proxy for $\kappa$ rather than sensing purely lubrication-specific signal changes directly.

This confounding does not invalidate the regression results: a model that uses speed-mediated amplitude changes to predict $\kappa$ is still predicting $\kappa$, and its predictions may be operationally useful. However, it does indicate that the model's generalisation to operating conditions outside its training distribution (different speed ranges, different bearing loads) may be limited. The two-stage analysis of Section~\ref{subsec:confound} makes this concrete: predicted operating points alone, fed through the ISO~281 mapping, recover $R^2 = \resTwoStageRsq$ of the direct model's $R^2 = \resHOrsqLgbAe$---the model is, to first order, an operating-point soft sensor, and its lubrication sensitivity beyond the operating point rests on the within-step evidence of Section~\ref{subsec:confound}. Future feature engineering should prioritise amplitude-invariant spectral-shape features (spectral skewness and kurtosis, spectral bandwidth, spectral flatness) that are expected to carry
a higher proportion of genuine lubrication-state information relative to operating
condition proxies.

\subsection{Implications for the $\kappa$ regression framework}
\label{subsec:disc_kappa}

The $\kappa$ regression framework is validated as a methodologically sound approach for passive non-invasive LCM, when $\kappa$ is assumed to be representative of the lubrication conditions in the bearing. A hold-out $R^2 = \resHOrsqLgbAe$ and RMSE~$= \resHOrmseLgbAe$~$\kappa$-units for a single AE sensor represents strong performance across a demanding continuous speed-and-temperature ramp spanning boundary through full-film lubrication conditions. Under matched random (sweep-level) splits this performance is on par with the passive $\kappa$ regression reported by Jakobsen et al.~\cite{jakobsen_detecting_2023,jakobsen_low_2024}. The figures reported here are not directly comparable to most prior work, however, because they are obtained under a stricter \emph{visit-grouped} evaluation protocol (Section~\ref{subsec:modelling_design}) that removes the near-duplicate leakage admitted by a sweep-level random split; to our knowledge they are the first $\kappa$-regression results reported under such a leakage-controlled protocol. The contribution is thus methodological---a more conservative, deployment-representative performance estimate---rather than a raw accuracy gain over earlier studies.

The repeated nested cross-validation design provides appropriately conservative
performance estimates: the \resNouterScores~outer RMSE scores capture fold-to-fold variance without the optimistic bias introduced by simple $k$-fold evaluation with shared hyperparameter search. The corrected significance tests confirm that performance differences between model families are assessed under assumptions that properly account for data re-use.

\subsection{Limitations}
\label{subsec:disc_limits}

\begin{enumerate}
  \item \textbf{Single bearing and lubricant.} All results are specific to the SYJ25 bearing and Keratech~22 oil. Generalisability to other bearing types, sizes, or lubricants is unknown.
  \item \textbf{Heterodyned US channel.} The passive US probe down-mixes a narrow $\approx$38--40~kHz sensing band to an audio-frequency baseband rather than recording wideband ultrasound directly, and its lubrication-relevant information is concentrated below $\approx$\SI{20}{kHz} (Section~\ref{subsec:origins}). It should not be equated with a wideband passive ultrasound sensor, and the US-channel results should be interpreted in light of this hardware constraint.
  \item \textbf{Incomplete operating-condition coverage.} The sampled operating points do not form a full speed--temperature factorial: high rotational speeds are reached predominantly at lower temperatures and high temperatures predominantly at lower speeds, leaving the high-speed/high-temperature corner sparsely populated (Figure~\ref{fig:op_conditions_distribution}). The $\kappa$ distribution is correspondingly skewed toward the boundary and mixed-lubrication regime, with full-film states ($\kappa > 1$) under-represented. Extrapolation into the under-sampled region of the operating space should therefore be treated with caution.
  \item \textbf{Model-derived $\kappa$ reference.} The viscosity ratio is not measured directly but computed for each sweep from the measured operating point---rotational speed and housing temperature---via the ISO~281 relations and the ASTM~D341 Walther viscosity--temperature equation (Section~\ref{subsec:kappa}). The regression target is therefore a model-derived reference that inherits the assumptions of those relations, notably the use of housing temperature as a proxy for the elastohydrodynamic contact temperature; this also underlies the operating-point soft-sensor interpretation of Section~\ref{subsec:disc_confound}.
\end{enumerate}
